\bookchapter{第三十章：投屏、演示與車載互聯}

\chapterintro{手機螢幕再大也就7英寸，但你的電視有55英寸、你的投影儀有100英寸。把手機畫面“扔”到大螢幕上，體驗是完全不同維度的。這一章講清楚投屏、演示和車載互聯這三個高頻場景，幫你把手機的能力延伸到更大的螢幕上。}

\booksection{一、無線投屏三大陣營——別搞混了}

無線投屏看起來都是“手機畫面到電視上”，但底層技術有三種，互不通用：

\begin{center}
\begin{tabularx}{\linewidth}{|X|X|X|}
\hline
\textbf{技術} & \textbf{誰用} & \textbf{特點} \tabularnewline
\hline
AirPlay & iPhone/iPad/Mac & 蘋果生態，體驗最流暢。需要Apple TV或支援AirPlay的電視 \tabularnewline
Miracast & 安卓手機 & Wi-Fi Direct直連，不需要路由器。安卓手機"無線投屏"預設用這個 \tabularnewline
DLNA & 所有平臺 & 只推送媒體檔案（影片/音樂/圖片），不映象螢幕。適合推影片到電視上播放 \tabularnewline
\hline
\end{tabularx}
\end{center}

\booksubsection{各平臺操作方法：}

\begin{itemize}
  \item \textbf{iPhone}：控制中心 > 螢幕映象 > 選擇你的電視/Apple TV。如果電視不支援AirPlay，可以買一個Apple TV盒子或支援AirPlay的電視棒。
  \item \textbf{安卓}：下拉通知欄 > 找到"無線投屏"/"智慧投屏"/"多屏互動"按鈕 > 選擇電視。需要手機和電視連同一個Wi-Fi。
  \item \textbf{萬能方案}：如果電視太老不支援任何無線投屏，買一個Chromecast或小米盒子/華為盒子插到電視HDMI口上，就能接收投屏了。
\end{itemize}

\booksection{二、有線投屏——更穩定、零延遲}

無線投屏方便但有延遲，玩遊戲和做演示時體驗不好。有線投屏能解決這個問題：

\begin{itemize}
  \item \textbf{iPhone}：需要一根Lightning/USB-C轉HDMI的轉換線（約100-200元）。插上後手機畫面直接輸出到電視/投影儀的HDMI口。
  \item \textbf{安卓}：部分手機支援USB-C轉HDMI直連（需要手機支援DP Alt Mode，大部分旗艦都支援）。買一根USB-C to HDMI線即可。
  \item \textbf{會議室場景}：如果你經常需要用手機做演示，建議在包裡常備一根轉HDMI線。會議室的電視/投影儀幾乎都有HDMI口，插上就能用，比折騰無線投屏靠譜多了。
\end{itemize}

\booksection{三、手機變電腦——桌面模式}

部分安卓手機支援“桌面模式”——連上顯示器後，手機會輸出一個類似電腦桌面的介面，可以用滑鼠鍵盤操作：

\begin{itemize}
  \item \textbf{Samsung DeX}：三星旗艦的招牌功能。連線顯示器後彈出桌面介面，可以多視窗操作，體驗最接近真正的電腦。配合藍芽鍵鼠，出差可以不帶筆記本。
  \item \textbf{華為桌面模式}：華為Mate系列/Pura系列連線顯示器後也有類似功能。支援多視窗、檔案管理、辦公App。
  \item \textbf{摩托羅拉Ready For}：Moto旗艦也支援，功能類似DeX。
  \item \textbf{適用場景}：輕辦公（回郵件、改檔案）、演示PPT、大屏看影片。不適合重度Excel操作或專業設計。
\end{itemize}

\booksection{四、車載互聯——開車用手機的正確方式}

開車時操作手機是違法的，但導航、聽音樂、接電話又是剛需。車載互聯絡統讓你在車機螢幕上使用手機的核心功能，手機可以放在一邊不碰：

\begin{itemize}
  \item \textbf{Apple CarPlay}：iPhone使用者的首選。透過USB或無線連線車機，車機螢幕上出現iPhone的精簡介面——導航（高德/百度）、音樂、電話、資訊、播客。Siri語音控制全程免手。
  \item \textbf{Android Auto}：安卓使用者對應方案。功能與CarPlay類似，支援Google Maps、Spotify等。國內因Google服務限制，部分功能不可用。
  \item \textbf{華為HiCar}：華為自研方案，國內車型支援率高。除了導航音樂，還能用華為視訊通話、泊車影像融合等特色功能。
  \item \textbf{百度CarLife}：跨平臺方案，iOS和安卓都能用。功能較少但相容性廣，很多老車型只支援CarLife。
\end{itemize}

\begin{tcbwarn}{安全提醒}
無論哪種車載互聯方案，開車時都應儘量使用語音控制。導航提前設好目的地，音樂提前選好歌單，不要在行駛中操作車機螢幕。

\end{tcbwarn}

\booksection{五、會議室投屏實戰——不再手忙腳亂}

開會是投屏最高頻的場景，也是最容易出問題的場景。以下是幾條實戰經驗：

\begin{itemize}
  \item \textbf{提前到場測試}：永遠不要等到會議開始才連投屏。提前10分鐘到場，測試好連線方式。Wi-Fi投屏失敗的機率比你想的高得多。
  \item \textbf{準備有線方案做Plan B}：在包裡放一根轉HDMI線。無線投屏連不上的時候，有線就是救命稻草。
  \item \textbf{關閉通知預覽}：投屏時你的微信訊息會顯示在會議室大屏上。投屏前在設定裡關閉"鎖屏通知內容預覽"，或者開啟"演示模式"（部分手機支援，投屏時自動隱藏通知）。
  \item \textbf{用"演示檔案"而不是"操作手機"}：如果你要展示PPT，直接把PPT存到手機裡用WPS開啟全屏播放，而不是投屏整個手機桌面。這樣更專業、更流暢。
\end{itemize}

\begin{tcbpractice}{動手：完成一次手機到電視的投屏}
\begin{enumerate}
  \item 確保手機和電視連同一個Wi-Fi。
  \item 用手機開啟一張照片或一個影片。
  \item 嘗試無線投屏（iPhone用螢幕映象，安卓用無線投屏按鈕）。
  \item 如果失敗，檢查電視是否支援對應協議，或嘗試DLNA推送（在相簿/影片App裡點"投屏"按鈕）。
  \item 投屏成功後，試試在手機上做其他操作，看看大屏是否同步顯示。
\end{enumerate}

\end{tcbpractice}

\begin{tcbsummary}{本章小結}
\begin{itemize}
  \item 無線投屏三陣營：AirPlay（蘋果）、Miracast（安卓）、DLNA（跨平臺推媒體）
  \item 有線投屏更穩定零延遲，會議室場景必備轉HDMI線
  \item Samsung DeX/華為桌面模式讓手機變身輕辦公電腦
  \item 車載互聯：CarPlay/Android Auto/HiCar/CarLife四大方案
  \item 會議室投屏鐵律：提前測試、準備Plan B、關閉通知預覽
\end{itemize}

\end{tcbsummary}
